Ein System ist ein LTI-System (Linear, Time Invariant), wenn für beliebige Eingangssignale $x_1(t), x_2(t)$ und Ausgänge $y_1(t), y_2(t)$ gilt:
\begin{description}
    \item[Linearität] $\mathcal{T}(ax_1(t) + bx_2(t)) = ay_1(t) + by_2(t)$
    \item[Zeitinvarianz] $\mathcal{T}(x(t-t_0))=y(t-t_0)$
\end{description}

\textbf{Eigenschaften}
\begin{description}
    \item[stabil] falls $\int_{-\infty}^\infty h(t)\diff t < \infty$
    \item[kausal] Wert des Ausgangssignals hängt nur vom Eingangssignal ab $\Leftrightarrow h(t)=0 \, \forall t<0$. Es gibt nicht kausale LTI-Systeme
\end{description}

\subsection{Graphische Faltung}
\begin{sectionbox}
    $t>0$ schiebt nach rechts! \\
    Faltung stückw. konst. Funktionen ergibt stück. lineare Funktionen $\Rightarrow$ Wert am Ende jedes Bereichs ausrechnen und mit Geraden verbinden
    \begin{enumerate}
        \item $h(\tau)$ zeichnen mit $\tau$-Achse
        \item Schnittpunkte von $h(\tau)$ mit der $\tau$-Achse mit dem $\tau$-Wert beschriften
        \item $h(-\tau)$ auf Glassichtfolie zeichnen. Horizontale Achse hat Ursprung bei $t$ anstatt $0$
        \item Für $t=0$ liegen die Schnittpunkte von $h(t-\tau)$ bei $\tau_i \Rightarrow$ für $t$ variabel liegen sie bei $t+\tau_i$. Diese an die Schnittpunkte auf der Glassichtfolie zeichnen
        \item Folie von links nach rechts über $h(\tau)$ ziehen. Jede Überschneidung einer sich bewegenden Kante (Folie) mit einer stat. Kante ($h(\tau)$) ergibt einen Bereich
        \item Bewegte Kante $=$ stat. Kante setzen und nach $t$ auflösen (z.B. $t+1s=-2s$)
        \item In Bereiche einteilen und für jeden Bereich eine feste und eine von $t$ abh. Grenze finden. Hierzu einfach den Integrationsbereich anschauen
    \end{enumerate}

    \textbf{Faltung mit Dirac $h(t) = \alpha\delta(t - T)$} \\
    $(h * x)(t)=\alpha x(t-T) \Rightarrow$ Ausgangssignal wird um $T$ nach rechts verschoben und mit $\alpha$ skaliert
\end{sectionbox}

\subsection{Bode-Diagramm}
\subsubsection{Logarithmische Rechnung}
\begin{description}
    \item[Maß] logarithmiertes positives Verhältnis zweier Leistungen jeweils gleicher Dimensionen, wobei die Bezugsgröße keinen festen Wert hat.
          z.B. $100W$ in den Lautsprecher hinein, $0.1W$ akustisches Signal heraus $\rightarrow \log_{10}(\frac{100W}{0.1W})=3dB$
\end{description}

\begin{sectionbox}
    \begin{emphbox}
        Das Dämpfungsmaß (wie stark ein Signal gedämpft wird) ist mit Eingangsleistung $P_1$ und Ausgangsleistung $P_2$ gegeben durch
        \begin{equation*}
            A=10\log_{10}\left(\frac{P_1}{P_2}\right)\text{dB}
        \end{equation*}

        Das Verstärkungsmaß $G$ ist das negative Dämpfungsmaß. Dieses wird z.B. bei Bode-Diagrammen aufgetragen
    \end{emphbox}

    Für Spannungen gilt $P=\dfrac{x_{eff}^2}{R}$ und somit \begin{equation*}
        L_V=20\log_{10}\left(\frac{x_{1eff}}{x_{0eff}}\right) \text{dB}
    \end{equation*}
\end{sectionbox}

\subsection{System erster Ordnung}
\subsection{System zweiter Ordnung}
TODO: Graph von h, s und Übertragungsfkt.
TODO: Normierung auf Q (Qualitätsfaktor)
